\section{Background and Problem Formulation}

\subsection{Reverse documentation engineering}

The original Reversa framework addresses a practical mismatch between agentic software maintenance and legacy codebases: coding agents need explicit behavioral contracts, while mature systems often encode critical rules only implicitly. Reversa structures discovery as a sequence of specialized roles that map project structure, excavate implementation details, infer domain and architecture, synthesize specifications, and review unresolved gaps \cite{macedo2026reversa}. Confidence is represented explicitly so that uncertain extractions are not silently treated as confirmed facts.

ReversaFeynman preserves this foundation. The extension does not redefine reverse documentation engineering; it adds a governance layer for what happens after claims, specifications, candidate repairs, trajectories, and learned policies begin to interact.

\subsection{Epistemic failure modes in agentic software engineering}

Repository-scale agents can fail even when individual model outputs appear plausible. We focus on six recurrent failure modes:

\begin{enumerate}
  \item \textbf{label substitution}: a named pattern or architecture is repeated without a causal explanation;
  \item \textbf{provenance loss}: a strong claim is detached from the code, test, execution, log, dataset, or artifact that supports it;
  \item \textbf{observation--inference collapse}: an interpretation is represented as if it had been directly observed;
  \item \textbf{non-falsifiable validation}: a test suite passes, but no meaningful oracle exists that could reveal a relevant behavioral error;
  \item \textbf{cargo-cult repair}: an established pattern is copied despite insufficient evidence that it addresses the local requirement;
  \item \textbf{confidence substitution}: memory, model confidence, human fluency, reward, or benchmark score is treated as equivalent to truth.
\end{enumerate}

The design of ReversaFeynman follows a Popperian intuition: strong engineering claims should expose conditions under which they could be shown false \cite{popper1959logic}. This does not turn software development into scientific experimentation in a strict philosophical sense; rather, it motivates operational gates that demand provenance, counterexamples, and executable checks when available.

\subsection{Epistemic states}

The framework uses four system-level states:

\begin{center}
\begin{tabular}{p{3.0cm}p{10.5cm}}
\toprule
\textbf{State} & \textbf{Meaning} \\
\midrule
\texttt{OBSERVED} & Supported by direct, traceable evidence accepted by the evidence governor. \\
\texttt{INFERRED} & Reasoned from available evidence but not directly established. \\
\texttt{UNVERIFIED} & Plausible or proposed, but not yet verified. \\
\texttt{BLOCKED} & Verification cannot proceed because required evidence, access, or a decision is missing. \\
\bottomrule
\end{tabular}
\end{center}

Human source quality is modeled separately from system evidence. A domain specialist may be classified as \texttt{HUMAN-VALIDATED}, \texttt{HUMAN-PARTIAL}, or \texttt{HUMAN-CONFLICT}; none of these labels is automatically equivalent to \texttt{OBSERVED}. This separation is essential when human explanation is the best available source but still requires traceability or executable corroboration.